\documentclass[11pt,a4paper]{article}
\usepackage[margin=25mm]{geometry}
\usepackage{amsmath,amssymb,amsthm,hyperref,microtype}
\setlength{\emergencystretch}{2em}
\newtheorem{theorem}{Theorem}
\newtheorem{lemma}[theorem]{Lemma}
\newtheorem{corollary}[theorem]{Corollary}
\newcommand{\GL}{\operatorname{GL}}
\newcommand{\SL}{\operatorname{SL}}
\newcommand{\PSL}{\operatorname{PSL}}
\newcommand{\PGL}{\operatorname{PGL}}
\newcommand{\diag}{\operatorname{diag}}
\newcommand{\F}{\mathbb F}
\title{Central quotients of the determinant-sign subgroup}
\author{Human Compression project: a component of the N012 recovery}
\date{}
\begin{document}\maketitle
For an odd prime power $q$ and $n\geq2$, put
\[
 H=\{g\in\GL_n(q):\det g\in\{1,-1\}\},\qquad
 Z_1=Z(H),\quad Z_2=Z(\SL_n(q)),\quad Q=H/Z_2.
\]
These are the groups denoted $\GL_n^{\pm1}(q)$ and
$\PGL_n^{\pm1}(q)$ in Markovskaya's thesis [T]. In particular, $Q$
is not defined by dividing by the full centre of $H$.
We first determine both quotients. We then reuse the four-dimensional
construction and fixed-space arguments from the reviewed recovery [L,U].
This proves specified parts of Chapter 4, including its complete
four-dimensional quotient case. It is not a recovery of the whole thesis.

\section{The two centres}
Write $d=\gcd(n,q-1)$ and $e=(q-1)/d$.
\begin{theorem}\label{centres}
The centres are
\[
 Z_1=\{\lambda I:\lambda^{2n}=1\},\qquad
 Z_2=\{\lambda I:\lambda^n=1\}.
\]
There are precisely two cases:
\begin{enumerate}
\item If $e$ is even, then $|Z_1:Z_2|=2$,
      $Q\cong\PSL_n(q)\times C_2$ and $H/Z_1\cong\PSL_n(q)$.
\item If $e$ is odd, then $Z_1=Z_2$ and the natural projective map
      embeds $Q$ as $\PSL_n(q)\rtimes C_2$ in $\PGL_n(q)$.
      The complement is represented by $\diag(-1,1,\ldots,1)$.
\end{enumerate}
The first case always holds when $n$ is odd; it can also hold when
$n$ is even.
\end{theorem}
\begin{proof}
The group $H$ contains every elementary transvection
$t_{ij}(a)=I+aE_{ij}$. A matrix commuting with all $E_{ij}$, $i\ne j$,
is scalar: the equation $gE_{ij}=E_{ij}g$ first makes its off-diagonal
entries zero and then equates every pair of diagonal entries.
Thus the centres in question consist precisely of the scalars satisfying
their determinant conditions. This gives the displayed formulae.

Choose a generator $u$ of $\F_q^*$. A scalar of determinant $-1$ exists
if and only if $nx\equiv(q-1)/2\pmod{q-1}$ is soluble, equivalently
$d\mid(q-1)/2$, or $e$ is even. In that case choose such a scalar $z$.
The determinant map shows that $H=\SL_n(q)\langle z\rangle$ and
that $\langle zZ_2\rangle$ is central of order two in $Q$, disjoint
from $\SL_n(q)/Z_2$. Hence $Q$ is the asserted direct product.
Dividing by its remaining central factor gives $H/Z_1\cong\PSL_n(q)$.

If no such scalar exists, all central scalars lie in $\SL_n(q)$, so
the projective map has precisely kernel $Z_1=Z_2$. Its image contains
$\PSL_n(q)$ with index two. The involution $D=\diag(-1,1,\ldots,1)$
has image outside $\PSL_n(q)$: otherwise some scalar multiple of $D$
would have determinant one, supplying a scalar of determinant $-1$.
Thus $D$ supplies a complement. Finally, for odd $n$ the scalar $-I$
already has determinant $-1$.
\end{proof}

For instance, when $n=6,q=5$, the scalar $2I$ has determinant
$2^6=-1$ in $\F_5$. Here $|Z_1|=4$ and $|Z_2|=2$;
$Q\cong\PSL_6(5)\times C_2$, whereas $H/Z_1\cong\PSL_6(5)$.
This corrects the assertion on [T, p.~69] that $Z_1=Z_2$ for every
even $n$, and the resulting universal index-two description in
Corollary 4.2.1. It does not refute the existence of the polytopes
asserted there. Their prescribed types require the separate marked
construction and are not proved merely by these group isomorphisms.

\section{Lifting a string marking through a central factor}
\begin{lemma}\label{lift}
Let $G$ be a finite perfect group with a rank-three string C-marking
$(a,b,c)$, so $ac=ca$. If $z$ generates $C_2$, then
$((a,z),(b,z),(c,z))$ is a string C-marking of $G\times C_2$.
\end{lemma}
\begin{proof}
Let $M$ be the subgroup generated by these three lifts. Projection
onto $G$ is onto. If its kernel were trivial, $M$ would be the graph
of a homomorphism $G\to C_2$. Such a homomorphism is trivial because
$G$ is perfect, whereas the three specified generators all have
second coordinate $z$. Therefore the kernel is $C_2$, and $M$ is
the full direct product.

For the end group $\langle a,b\rangle$, put $m=|ab|$.
The standard dihedral presentation
$\langle a,b\mid a^2=b^2=(ab)^m=1\rangle$ has a homomorphism
to $C_2$ sending both reflections to $z$. Consequently
$\langle(a,z),(b,z)\rangle$ is its graph and projects injectively
onto $\langle a,b\rangle$. The same argument applies to
$\langle(b,z),(c,z)\rangle$. Their intersection projects into
$\langle b\rangle$ by the given intersection property and is
therefore exactly $\langle(b,z)\rangle$. Each end pair is itself
a rank-two C-marking. Together with the distant commutation and
distinct nonidentity generators, this verifies the rank-three
intersection property.
\end{proof}

We retain the classification of admissible simple linear groups and
the Conder--Oliveros criterion already explicitly acquired in [L].
Here admissible means generated by three elements of order at most two,
two of which commute. For $n\geq4$ and odd $q$, the simple group
$\PSL_n(q)$ is admissible. The criterion turns an admissible triple
of distinct involutions in a finite group with no nontrivial cyclic
normal subgroup into a rank-three string C-marking. These are external
classification and geometric inputs, not new theorems of this note.

\begin{corollary}\label{centralexistence}
For $n\geq4$, odd $q$ and even $e$, the group $Q$ is the automorphism
group of a regular abstract polytope of rank three.
\end{corollary}
\begin{proof}
The retained simple-group result gives a string C-marking of
$\PSL_n(q)$. This group is nonabelian simple, hence perfect.
Apply Lemma~\ref{lift} and Theorem~\ref{centres}, followed by the
standard group--polytope correspondence.
\end{proof}
This proves existence, without claiming the types $[4p,n]$ or
$[4p,2n]$ of [T, Corollaries 4.2.1--4.2.3].

\section{The four-dimensional quotient for every odd field}
The constructions in [L, Section 4] were used there under
$\gcd(4,q-1)=2$ to obtain all of $\PGL_4(q)$. We isolate the preceding
generation argument, which needs only that $q$ is odd.
\begin{lemma}\label{four}
For every odd $q$, the group $H_4(q)$ has generators $A,B,G$ such
that $A^2,B^2,G^2\in\{I,-I\}$, $AB=-BA$, and
$\det A=\det B=1$, $\det G=-1$.
\end{lemma}
\begin{proof}
Choose a primitive element $u\in\F_q^*$ and take
\[
 A=\begin{pmatrix}0&0&0&-1\\0&0&-1&0\\0&1&0&0\\1&0&0&0\end{pmatrix}.
\]
In characteristic three use
\[
 B=\begin{pmatrix}0&0&-1&0\\0&0&0&1\\1&0&0&0\\0&-1&0&0\end{pmatrix},
 \qquad G=\begin{pmatrix}-1&0&0&0\\-1&1&0&0\\-u&0&1&0\\-u&0&0&1\end{pmatrix}.
\]
Use exactly the words $J,D,T,P$ and their identities in
[L, Section 4.1]. Those are polynomial identities over $\F_3[u]$.
The only nonvanishing condition used before obtaining $\SL_4(q)$
is $w=u^3+u\ne0$. It holds for every primitive $u$ in characteristic
three: otherwise $u^2=-1$, so $u$ has order four and $q-1=4$,
contrary to $q$ being a power of three.

The displayed roots in that section give a strongly connected graph
on coordinates $1,2,3$. The root-parameter theorem of [U] supplies
its diagonal sign matrices. In particular, isolating the factors of
$T=t_{24}(u)t_{34}(u^2)$ gives $t_{24}(-u)$ and its conjugate
$t_{42}(u)$. The monomial permutations generated on the first three
coordinates and by $AB$ give all four coordinates. Hence all roots
are reached. The opposite-root product $-u^2$ makes the parameter
field equal to $\F_q$, by the primitive-square test in [U].
This is the complete generation argument of [L] with its field
condition checked; its later assertion about scalar multiples of $A$
generating $\GL_4(q)$ is not used.

In odd characteristic different from three use instead
\[
 B=\begin{pmatrix}0&0&1&0\\0&0&0&1\\1&0&0&0\\0&1&0&0\end{pmatrix},
 \qquad G=\begin{pmatrix}-1&0&0&0\\1&1&0&0\\u&0&1&0\\u&0&0&1\end{pmatrix}.
\]
The identities in [L, Section 4.2] give the two ordinary swaps
$(12),(34)$, the matrix $\diag(-1,-1,1,1)$, and the root
$(UE_0)^2=t_{24}(3u)$, with $U,E_0$ defined there. They involve
only division by two and three in the prime field. Thus they give
$t_{24}(u)$ and $t_{42}(u)$ for every field now considered.
The monomial matrices reach every ordered pair across
$\{1,2\}\mid\{3,4\}$; this directed graph is strongly connected.
Root propagation and the opposite parameter product $u^2$ again
give $\SL_4(q)$. No nonsquare assumption on $-1$ enters this argument.

In both cases the square, commutation and determinant assertions
follow from the displayed matrices. All generators lie in $H$,
they contain $\SL_4(q)$, and $G$ has determinant $-1$.
Consequently they generate $H$, as claimed.
\end{proof}

\begin{theorem}\label{qfour}
For every odd prime power $q$, $Q_4(q)$ is the automorphism group
of a regular abstract polytope of rank three.
\end{theorem}
\begin{proof}
Since $-I\in Z(\SL_4(q))$, the images of the three matrices of
Lemma~\ref{four} are involutions, with the first two commuting.
They are distinct and nonidentity: none of the displayed matrices
is scalar, $AB$ is not scalar, and the image of $G$ has nontrivial
determinant coset while those of $A,B$ do not. They generate $Q$.

If $e$ is even, Corollary~\ref{centralexistence} applies. We make
no assertion here that this particular matrix triple has the
intersection property in that case; the central-lift argument
supplies a marking.
If $e$ is odd, Theorem~\ref{centres} embeds $Q$ between
$\PSL_4(q)$ and $\PGL_4(q)$. The centralizer of $\PSL_4(q)$ in
$\PGL_4(q)$ is trivial. Indeed, a projective centralizer conjugates
each transvection to a scalar multiple of itself. Its sole eigenvalue
is one, forcing that scalar to be one; centralizing all matrix units
then forces a scalar matrix. Any cyclic normal subgroup of $Q$
intersects the nonabelian simple socle trivially and therefore
centralizes it, since the commutator lies in the intersection.
It too is trivial. The retained Conder--Oliveros criterion applies
to the displayed triple.
\end{proof}

The printed matrix $\beta$ on [T, p.~72] instead satisfies
$\beta^2=\diag(-1,-1,1,1)$, which is not scalar in odd characteristic.
Its image in $Q$ is not an involution. This defect was already recorded
in the project's phase 22 source audit; it is not a new discovery here.
Lemma~\ref{four} replaces that witness using the earlier journal
construction. Theorem~\ref{qfour} recovers the intended existence
statement without relying on the defective matrix.

\section{The small linear groups}
\begin{theorem}\label{small}
For odd $q$ and $2\leq n\leq4$, the group $H_n(q)$ is not admissible.
\end{theorem}
\begin{proof}
The irreducibility proof for $\operatorname{Sym}^2(\F_q^n)$ in [U]
uses only diagonal sign matrices, permutation matrices and elementary
transvections. All belong to $H$, so that proof establishes
irreducibility for $H$ as well. If $a,b,c$ were an admissible
triple with $ab=ba$, the five factors $a,b,ab,c,c$ would have product
one and generate this irreducible representation. Scott's inequality,
proved in [U], therefore gives
\[
 \sum_{g\in(a,b,ab,c,c)}\operatorname{codim}\operatorname{Fix}(g)
 \ \geq\ n(n+1).
\]
For an involution with negative eigenspace dimension $r$ on the
natural module, its symmetric-square fixed-space codimension is
$r(n-r)$: the negative part consists of the products of one positive
and one negative eigenvector. This also covers the identity.
For $n=2$ the sum is at most five, less than six. For $n=3$ it is
at most ten, less than twelve. For $n=4$ it is at most twenty,
so equality is forced in all five terms. In particular $a,b,c$
each have two negative eigenvalues and determinant one. They cannot
generate $H$, which contains $\diag(-1,1,1,1)$.
\end{proof}

\section{Attribution and remaining work}
The theorem statements being recovered are due to Markovskaya [T].
The four-dimensional matrices belong to the Markovskaya--Nuzhin
journal construction recovered in [L]; the fixed-space argument is
the source's Scott and symmetric-square method. The root-parameter
proofs, corrected matrix identities and primitive-square test are
acquired from [U,L]. The present contribution is the complete
centre calculation in arbitrary dimension, the explicit accounting
of both quotients, the central-lift argument, and the proof that the
existing four-dimensional construction works for every odd field.
These arguments are not asserted to be new in the literature.

The special four-dimensional centre calculation and the defective
printed witness were already recorded in phase 22 and are reused.
The complete proofs in [U,L] are dependencies; they must be included
in any later cost comparison wherever not already acquired. This
note does not use the retained integral generation theorem of [L].

Outstanding for N012 are the integral constructions of Chapter 2,
the remaining uniform constructions for $H_n(q)$ in Chapter 4,
the complete two-dimensional quotient cases, the noncentral
even-dimensional quotient cases above dimension four, and the
claimed polytope types. Chapter 3 is mapped to the reviewed N011
recovery, subject to a complete source-by-source routing audit.
No whole-thesis completion, open-question resolution, Lean
certificate, or compression ratio is claimed here.

\begin{thebibliography}{9}\small
\bibitem{T}[T] I.~A.~Markovskaya, \emph{Linear groups over finite fields
and over the ring of integers as automorphism groups of regular
polytopes of rank 3}, thesis, Siberian Federal University, 2025,
80-page supplied edition; especially pp.~53--75.
\bibitem{L}[L] \emph{Three involutions in finite general linear groups},
reviewed N011 recovery of I.~A.~Markovskaya and Ya.~N.~Nuzhin,
\emph{On generation of the groups $\GL_n(q)$ and $\PGL_n(q)$ by three
involutions, two of which commute}, Trudy Inst. Mat. Mekh. UrO RAN
31(4) (2025), 247--259. DOI 10.21538/0134-4889-2025-31-4-fon-03.
\bibitem{U}[U] \emph{Root parameters and fixed-space obstructions},
reviewed common paper accompanying [L]. Its proofs specialize
classical results of Levchuk and Scott, as attributed there.
\end{thebibliography}
\end{document}
